%---------------------- Problema 1 ------------------------------
\section*{Problema 1}

\begin{enunciado}{Enunciado}
    % TODO: Pegar aqu\'i el enunciado del Problema 1 de la Tarea 2.
    Reemplazar este texto por el enunciado completo del problema
    (datos, geometr\'ia, propiedades del material, condiciones de borde,
    y lo que se pide calcular/discutir).
\end{enunciado}

\vspace{0.7cm}

% Figura del problema (opcional). Ejemplo de plantilla TikZ, ajustar o
% eliminar seg\'un corresponda.
% \begin{center}
% \begin{tikzpicture}[scale=0.9, >=stealth, thick]
%
% \end{tikzpicture}
% \captionof{figure}{Descripci\'on de la figura.}
% \label{fig:t2-p1}
% \end{center}

%---------------------- Desarrollo -------------------------------

\begin{desarrollo}{a) Desarrollo te\'orico / planteamiento}
    % TODO: desarrollo anal\'itico o planteamiento del problema.
\end{desarrollo}

\begin{desarrollo}{b) Soluci\'on por MEF}
    % TODO: implementaci\'on, matrices de rigidez, resultados num\'ericos.
    % Ejemplo de bloque de c\'odigo Python con el mismo estilo usado en T1:
    % \begin{lstlisting}[style=pythonstyle, caption={Descripci\'on}, label={lst:t2-p1}]
    % def ejemplo():
    %     pass
    % \end{lstlisting}
\end{desarrollo}

\begin{desarrollo}{c) Comparaci\'on y discusi\'on}
    % TODO: comparar resultados, discutir convergencia, errores, etc.
\end{desarrollo}
